\chapter{Phase 2: ACI and Tool Engineering (The Hands Phase)}

The goal of this phase is to build robust interfaces that agents can reliably use. We move from simple functions to tools with built-in safety guardrails and error-handling loops.

\section{Step 4: Safe SQL Tool}
This step demonstrates building a read-only SQL tool. We implement essential ACI guardrails: enforcing \texttt{SELECT}-only queries, injecting \texttt{LIMIT} clauses, and catching SQL errors to return them as structured JSON for self-correction.

\begin{lstlisting}[caption={step4\_sql\_tool.py}, label={lst:step4}]
"""
PHASE 2 - STEP 4: ACI - Safe SQL Read Tool
Harness Pillar: ACI (Agent-Computer Interface)
Goal: Build a read-only SQL tool with safety guardrails.
"""

import sqlite3
import re
import json

MAX_ROWS = 50
FORBIDDEN_PATTERNS = [r"\bDROP\b", r"\bDELETE\b", r"\bINSERT\b", r"\bUPDATE\b"]

def safe_sql_query(query: str, db_path: str) -> dict:
    query = query.strip()
    # Guardrail: Only SELECT allowed
    if not query.upper().lstrip().startswith("SELECT"):
        return {"success": False, "error": "FORBIDDEN: Only SELECT allowed."}

    # Guardrail: Block mutation keywords
    for pattern in FORBIDDEN_PATTERNS:
        if re.search(pattern, query, re.IGNORECASE):
            return {"success": False, "error": f"FORBIDDEN: Disallowed keyword."}

    # Guardrail: Inject LIMIT
    if "LIMIT" not in query.upper():
        query = f"{query.rstrip(';')} LIMIT {MAX_ROWS};"

    try:
        conn = sqlite3.connect(db_path)
        conn.row_factory = sqlite3.Row
        cur = conn.cursor()
        cur.execute(query)
        rows = [dict(row) for row in cur.fetchall()]
        conn.close()
        return {"success": True, "rows": rows}
    except sqlite3.Error as e:
        return {"success": False, "error": str(e), "hint": "Check column names."}
\end{lstlisting}

\subsection*{Key Takeaways}
\begin{itemize}
    \item The ACI Observation Layer should be read-only by default to protect your data.
    \item Always return errors as structured JSON, not exceptions; the agent reads the error and self-corrects.
    \item Inject \texttt{LIMIT} automatically to prevent performance issues.
\end{itemize}

\section{Step 6: Error-Handling Refinement}
The refinement loop is what separates a fragile script from a resilient agent. By setting \texttt{is\_error=True} and providing the full stack trace, we enable Claude to fix its own mistakes autonomously.

\begin{lstlisting}[caption={step6\_error\_refinement.py}, label={lst:step6}]
"""
PHASE 2 - STEP 6: Error-Handling Refinement Loop
Harness Pillar: Standard Workflows & Refinement Loops
Goal: Feed tool errors back to the model so it can self-correct.
"""

import anthropic
import json

def run_self_correcting_agent(task: str):
    messages = [{"role": "user", "content": task}]
    while True:
        response = client.messages.create(
            model="claude-opus-4-5-20251101",
            max_tokens=2048,
            tools=TOOLS,
            messages=messages,
        )
        messages.append({"role": "assistant", "content": response.content})

        if response.stop_reason == "end_turn":
            return response.content[-1].text

        if response.stop_reason == "tool_use":
            tool_results = []
            for block in response.content:
                if block.type == "tool_use":
                    # Simulated execution with error
                    result = "NameError: name 'fib' is not defined"
                    is_error = True
                    
                    tool_results.append({
                        "type": "tool_result",
                        "tool_use_id": block.id,
                        "content": result,
                        "is_error": is_error,  # Critical signal for self-correction
                    })
            messages.append({"role": "user", "content": tool_results})
\end{lstlisting}

\subsection*{Key Takeaways}
\begin{itemize}
    \item \texttt{is\_error=True} in the tool result is the signal that activates self-correction.
    \item Include the full error or stack trace in the tool result content.
    \item The refinement loop is essential for production agents operating in real-world, error-prone environments.
\end{itemize}
